\documentclass[10.5pt,a4paper]{article}
% comm_scripts user guide: what every part of the repository does and how to run it.
% Build: tectonic comm_scripts_usage.tex   (or pdflatex twice).  Overleaf: upload comm_scripts_usage_overleaf.zip
\usepackage[margin=2.2cm]{geometry}
\usepackage[T1]{fontenc}
\usepackage{lmodern}
\usepackage{microtype}
\usepackage{xcolor}
\usepackage{booktabs}
\usepackage{tabularx}
\usepackage{longtable}
\usepackage{array}
\usepackage{enumitem}
\usepackage{listings}
\usepackage{tcolorbox}
\usepackage{titlesec}
\usepackage{fancyhdr}
% url before hyperref: keep spaces in code tokens and allow line breaks at spaces, / - _ .
\usepackage[obeyspaces,spaces]{url}
\usepackage[hidelinks]{hyperref}
\makeatletter\g@addto@macro{\UrlBreaks}{\do\_\do\.\do\=}\makeatother
\urlstyle{tt}

\definecolor{accent}{HTML}{1F4E79}
\definecolor{codebg}{HTML}{F4F6F8}
\definecolor{notebg}{HTML}{EAF2FB}
\definecolor{warnbg}{HTML}{FDF1E3}
\hypersetup{colorlinks=true, linkcolor=accent, urlcolor=accent, citecolor=accent}
\titleformat{\section}{\Large\bfseries\color{accent}}{\thesection}{0.8em}{}
\titleformat{\subsection}{\large\bfseries\color{accent}}{\thesubsection}{0.8em}{}
\titleformat{\subsubsection}{\normalsize\bfseries}{\thesubsubsection}{0.8em}{}
\setlength{\parindent}{0pt}
\setlength{\emergencystretch}{2.5em}
\setlength{\parskip}{0.55em}
\setlist{itemsep=0.15em, topsep=0.3em}
\renewcommand{\arraystretch}{1.18}
\pagestyle{fancy}
\fancyhf{}
\fancyhead[L]{\small\color{accent} comm\_scripts user guide}
\fancyhead[R]{\small\nouppercase{\leftmark}}
\fancyfoot[C]{\small\thepage}

\lstset{basicstyle=\ttfamily\small, backgroundcolor=\color{codebg}, frame=single, rulecolor=\color{codebg},
        breaklines=true, columns=fullflexible, keepspaces=true, xleftmargin=4pt, xrightmargin=4pt,
        aboveskip=0.6em, belowskip=0.6em, showstringspaces=false}
\lstdefinestyle{py}{language=Python, keywordstyle=\color{accent}\bfseries, commentstyle=\color{gray}}

% code tokens: \cc{...} is \path from the url package (_ -- ~ $ verbatim, spaces kept); not in moving arguments
\newcommand{\cc}{\path}
\newtcolorbox{note}[1][Note]{colback=notebg, colframe=accent, boxrule=0.5pt, arc=2pt, left=6pt, right=6pt,
                              top=4pt, bottom=4pt, fonttitle=\bfseries\small, title=#1}
\newtcolorbox{warn}[1][Caution]{colback=warnbg, colframe=orange!70!black, boxrule=0.5pt, arc=2pt, left=6pt,
                              right=6pt, top=4pt, bottom=4pt, fonttitle=\bfseries\small, title=#1}
\setlength{\LTcapwidth}{\textwidth}
\newcolumntype{L}{>{\raggedright\arraybackslash}X}
\newcolumntype{P}[1]{>{\raggedright\arraybackslash}p{#1}}

\title{\color{accent}\textbf{comm\_scripts: user guide}\\[0.4em]
       \large What the code does and how to run each part}
\author{skipper-CCD module testing}
\date{version of 2026-10-25 (branch \texttt{clustering-methods-2026-10-25})}

\begin{document}
\maketitle
\thispagestyle{empty}

\begin{abstract}
\noindent \texttt{comm\_scripts} takes raw skipper-CCD images to calibrated results with panaSKImg (WADERs). It
writes the panaSKImg configuration itself and measures gain, noise and dark current per CCD with errors. It can
correct the gain's dependence on pixel charge, run the Fe-55 composite analysis, and cluster with A\_GMX, a
clusterer that writes files in the WADERs format. The repository also contains standalone utilities, the research
code of the clustering study, and the scripts behind the analysis notes. This guide lists every part, shows how to
run it, and says which inputs it needs and which outputs it writes. Section~\ref{sec:recipes} covers the common
tasks in a few lines each. The physics behind the pipeline steps is in \cc{docs/PIPELINE_GUIDE.md}, and the A\_GMX
algorithm is described in \cc{docs/aGMX/aGMX_clustering.tex}.
\end{abstract}

\tableofcontents
\newpage

%==================================================================================================================
\section{Overview}
%==================================================================================================================
\subsection{What is in the repository}

\begin{center}
\begin{tabularx}{\textwidth}{@{}P{3.3cm}P{3.6cm}L@{}}
\toprule
\textbf{folder} & \textbf{kind} & \textbf{contents} \\
\midrule
\cc{bin/} & commands & \cc{run_WADERs} (alias \cc{runWADERs}): the full pipeline; \cc{run_AGMX}: A\_GMX
  clusters for a finished run \\
\cc{src/} & pipeline modules & \cc{autoconfig_panaSKImg.py}, \cc{gain_dc_split.py}, \cc{fit_quality.py},
  \cc{gain_vs_charge.py}, \cc{agmx_root.py} (all also runnable on their own) \\
\cc{utils/} & standalone helpers & \cc{compare_runs.py}, \cc{split_clustersRec_by_ext.py},
  \cc{compare_dcfit_vs_skips.py}, \cc{skip_noise_L1.py} \\
\cc{clustering_tests/} & research code & the clustering methods (A\_GMX, GMX, MHC, MFD+G, hybrid, \dots) and
  the scripts that tested them \\
\cc{tests/} & unit tests & simulated inputs, no data needed; run in under a minute \\
\cc{docs/} & guides, notes, reports & \cc{README.md} (index), this guide, \cc{PIPELINE_GUIDE.md}, the A\_GMX
  document, the Am-241 note, the gain-correction note and reports, \cc{notes/} (hand-over notes, open items) \\
\bottomrule
\end{tabularx}
\end{center}

Keep \cc{bin/} and \cc{src/} next to each other: \cc{run_WADERs} and \cc{run_AGMX} load their modules from
\cc{../src}. \cc{src/agmx_root.py} also imports the clustering code from \cc{../clustering_tests}.

\subsection{Requirements and set-up}
\begin{itemize}
  \item \textbf{Python environment}: numpy, scipy, astropy, matplotlib, PyROOT and pysimdamicm (panaSKImg). On
        this machine that is the conda environment \cc{damic}.
  \item \textbf{Automatic environment}: \cc{run_WADERs} and \cc{run_AGMX} can be started from any python. If
        that python lacks ROOT or pysimdamicm, they look for a conda/mamba environment that has panaSKImg and
        re-run themselves there with \cc{conda run}. Nothing is left activated afterwards.
  \item \textbf{Choosing the environment}: \cc{--env NAME} or \cc{$WADERS_ENV}. To choose the panaSKImg script,
        use \cc{--panaskimg PATH} or \cc{$PANASKIMG}.
  \item \textbf{Optional}: ROOT's \cc{root} executable, for \cc{--composite} (it runs the
        \cc{MakeComposite_Fe55_Clusters_Extensions.cc} macro). \cc{tectonic} or pdflatex, to build the
        LaTeX documents.
  \item \textbf{Calling by name}: to run the commands from anywhere, add
        \cc{export PATH="/path/to/comm_scripts/bin:$PATH"} to \cc{~/.zshrc}.
\end{itemize}

\begin{note}[Environment variables]
\begin{tabular}{@{}ll@{}}
\cc{WADERS_ENV} & conda environment to use (\cc{run_WADERs}, \cc{run_AGMX}) \\
\cc{PANASKIMG} & panaSKImg script \\
\cc{MAKECOMPOSITE_MACRO} & \cc{MakeComposite_Fe55_Clusters_Extensions.cc} used by \cc{--composite} \\
\cc{AM241_RUN_DIR} & data folder of the Am-241 note scripts (Section~\ref{sec:notes}) \\
\end{tabular}
\end{note}

\subsection{Conventions used throughout}
\begin{itemize}
  \item \textbf{Inputs}: a \emph{raw image} is a skip-averaged FITS image (\cc{.fz}, \cc{.fits}, \cc{.fits.bz2}),
        with one HDU per amplifier, i.e.\ per CCD (\cc{EXT1}--\cc{EXT4}). A \emph{list} is a \cc{.txt} file with
        one image path per line. Relative paths are relative to the \cc{.txt}; empty lines and \texttt{\#} lines are
        ignored.
  \item \textbf{Output folder}: every run writes \cc{WADERs_analysis_<run name>/}. The run name is the image name,
        or the \cc{.txt} name for a list.
  \item \textbf{Cluster files}: \cc{panaSKImg_clustersRec_<image>_ext<N>.root}, one per CCD, with trees
        \cc{clustersRec} (one entry per image, one vector element per cluster), \cc{info} and
        \cc{process_config}. Energies are in keV with 3.74\,eV per electron.
  \item \textbf{Modules}: the ACM board id appears in the image name
        (\cc{..._106_20260330_123052_21}: ACM 106). Images with the same ACM share the CCDs and are grouped,
        e.g.\ for one gain curve per module.
  \item \textbf{Fe-55}: results are quoted for front-side clusters ($\sigma_{xy}<1$\,px) unless stated otherwise.
\end{itemize}

%==================================================================================================================
\section{Common tasks in a few lines}\label{sec:recipes}
%==================================================================================================================
\begin{longtable}{@{}P{5.4cm}P{10.6cm}@{}}
\toprule
\textbf{I want to \dots} & \textbf{run} \\
\midrule
\endhead
analyse one image & \cc{run_WADERs image.fz} \\
analyse a list, 8 at a time, elsewhere & \cc{run_WADERs images.txt -o ~/analysis -j 8} \\
only gain, noise, dark current (fast) & \cc{run_WADERs images.txt --no-clusters} \\
gain-corrected cluster energies & \cc{run_WADERs images.txt --gain-correction} \\
compare per-image and per-module $g(n)$ & \cc{run_WADERs images.txt --gain-correction both} \\
Fe-55 composite plots per module & \cc{run_WADERs images.txt --gain-correction --composite --module 109=DM02} \\
cluster with A\_GMX in the pipeline & \cc{run_WADERs images.txt --gain-correction --clusterer agmx --composite} \\
WADERs and A\_GMX side by side & \cc{run_WADERs images.txt --gain-correction --clusterer both --composite} \\
A\_GMX for a run already done & \cc{run_AGMX WADERs_analysis_<run>} \\
cluster thresholds in electrons & \cc{run_WADERs images.txt --autoconfig-args "--cluster-threshold-e auto"} \\
check a code change changes nothing & run old and new code with \cc{-o old} / \cc{-o new}, then
  \cc{python utils/compare_runs.py old new} \\
one JSON for panaSKImg, by hand & \cc{python src/autoconfig_panaSKImg.py image.fz -o image.json} \\
split an \cc{_ext1234.root} file per CCD & \cc{python utils/split_clustersRec_by_ext.py f_ext1234.root -o per_ext/} \\
noise vs number of skips (level-1 data) & \cc{python utils/skip_noise_L1.py img_..._L1_....RJ45ADCDATA0 -o out/} \\
run the unit tests & \cc{python -m unittest discover -s tests -v} (repository root, \cc{damic} python) \\
\bottomrule
\end{longtable}

%==================================================================================================================
\section{The pipeline: \texttt{run\_WADERs}}\label{sec:runwaders}
%==================================================================================================================
\begin{lstlisting}
run_WADERs <image.fz | image.fits | image.fits.bz2 | list.txt> [-o PARENT] [-j N] [options]
run_WADERs -h          # the options
run_WADERs -h -v       # the full guide: output folders, what each step does, all autoconfig settings
\end{lstlisting}

\subsection{What it does, step by step}
Each step is one function in \cc{bin/run_WADERs}; \cc{main()} calls them in this order.

\begin{center}
\begin{tabularx}{\textwidth}{@{}P{4.3cm}L@{}}
\toprule
\textbf{step} & \textbf{what happens} \\
\midrule
\cc{check_request} & inputs exist, the output folder is free (else \cc{--overwrite}), the options are
  compatible \\
\cc{ensure_environment} & re-runs inside the conda environment with panaSKImg if needed \\
\cc{step1_configure} & per CCD, \cc{gain_dc_split.py --per-ext} writes two configurations, both measured from the
  raw pixels: a \emph{dark-current} JSON, which fits only the Poisson-like 0\,e$^-$/1\,e$^-$ peaks, and a
  \emph{gain} JSON, which fits all peaks. Both are fitted with panaSKImg. The final JSON holds the all-peaks gain
  if that is more precise (errors scaled by $\sqrt{\chi^2/\mathrm{ndf}}$) and agrees within $3\sigma$. Flooded
  or unresolved CCDs are flagged. \\
\cc{step2_run_panaskimg} & panaSKImg once per CCD (\cc{--ext N}), in parallel (\cc{-j}) \\
\cc{step3_collect_root_files} & cluster files go to \cc{root/}, dark-current fits to \cc{root/dc_fits/}, and the
  file list to \cc{root/root_files.txt} \\
\cc{step3b_gain_correction} & with \cc{--gain-correction}: measure $g(n)$ and write corrected cluster files
  (Section~\ref{sec:gain}) \\
\cc{step3d_agmx} & with \cc{--clusterer agmx|both}: A\_GMX cluster files in the same format
  (Section~\ref{sec:agmx}) \\
\cc{step3c_composite} & with \cc{--composite}: the MakeComposite macro per image type and module
  (Section~\ref{sec:composite}) \\
\cc{step4_summary} & \cc{fit_quality.csv} and \cc{summary.txt}: gain, noise and dark current per CCD with errors,
  which gain was used and why, and the results of the optional steps \\
\bottomrule
\end{tabularx}
\end{center}

\subsection{Options}
\begin{longtable}{@{}P{5.2cm}P{10.8cm}@{}}
\toprule
\textbf{option} & \textbf{meaning} \\
\midrule
\endhead
\multicolumn{2}{@{}l}{\textit{general}} \\
\cc{-o DIR}, \cc{--parent DIR} & folder in which \cc{WADERs_analysis_<run>} is created (default: current folder) \\
\cc{-j N}, \cc{--jobs N} & panaSKImg runs at the same time (default 4, or fewer CPUs) \\
\cc{--overwrite} & replace an existing output folder \\
\cc{--all-ext} & also run CCDs flagged as flooded/unresolved (slow; fits not reliable) \\
\cc{--no-clusters} & noise, gain, dark current only: no ClusterFinder (much faster) \\
\cc{--tune} & per CCD, choose the pedestal method and fit window with panaSKImg by the smallest total fit error
  ($\sim$1 min more per image) \\
\cc{--keep-work} & keep the intermediate fit runs in \cc{work/} \\
\cc{--autoconfig-args "..."} & settings for the JSON generation, in one quoted string (Table~\ref{tab:autoconfig}) \\
\cc{--env NAME}, \cc{--panaskimg PATH} & environment and panaSKImg script \\
\multicolumn{2}{@{}l}{\textit{gain versus charge} (Section~\ref{sec:gain})} \\
\cc{--gain-correction [MODE]} & \cc{pooled} (default: one curve per module), \cc{per-image}, \cc{both} \\
\cc{--gain-model} & \cc{interpolated} (default), \cc{broken}, \cc{linear} \\
\cc{--gain-combine} & \cc{coherent} (default), \cc{incoherent} \\
\cc{--gain-anchor} & \cc{comb} (default), \cc{fit} \\
\multicolumn{2}{@{}l}{\textit{clustering} (Section~\ref{sec:agmx})} \\
\cc{--clusterer} & \cc{waders} (default: panaSKImg's ClusterFinder only), \cc{agmx} (also A\_GMX files, used
  by \cc{--composite}), \cc{both} (\cc{--composite} on both) \\
\multicolumn{2}{@{}l}{\textit{composite} (Section~\ref{sec:composite})} \\
\cc{--composite} & run MakeComposite on the cluster files \\
\cc{--composite-macro PATH} & the macro (default \cc{$MAKECOMPOSITE_MACRO}, else the one in
  \cc{~/Documents/LSM_Module_Testing_Analysis/src/}) \\
\cc{--module ACM=NAME} & module name of an ACM, e.g.\ \cc{--module 109=DM02} (repeatable) \\
\cc{--image-number N} & image type for the macro (default: \cc{Image_<N>} in the name) \\
\bottomrule
\end{longtable}

\subsection{Output folder}
\begin{longtable}{@{}P{5.4cm}P{10.6cm}@{}}
\toprule
\textbf{path (in the run folder)} & \textbf{contents} \\
\midrule
\endhead
\cc{*.pdf}, \cc{*_EXT<n>_waders.fits.bz2} & panaSKImg plots; the calibrated image of each CCD (electrons) \\
\cc{summary.txt}, \cc{fit_quality.csv} & gain, noise, dark current per CCD with scaled errors and the
  dark-current bias \\
\cc{root/} & panaSKImg cluster files, one per CCD, and \cc{root_files.txt} \\
\cc{root/dc_fits/} & dark-current fits (histogram + function) as ROOT \\
\cc{json/<image>/final/} & the JSONs that were run (one per CCD) \\
\cc{json/<image>/gain/}, \cc{dc/} & intermediate JSONs; \cc{<image>_gain_dc_summary.txt} compares the fits \\
\cc{logs/} & panaSKImg, autoconfig and MakeComposite logs \\
\cc{gain_vs_charge/} & $g(n)$ curves (JSON) and their plots \\
\cc{root_gaincorr/}, \cc{root_gaincorr_per_image/} & gain-corrected cluster files \\
\cc{root_agmx/}, \cc{root_agmx_gaincorr*/} & A\_GMX cluster files \\
\cc{composite*/} & MakeComposite plots and the file lists given to it \\
\bottomrule
\end{longtable}

\begin{note}[Reading \cc{summary.txt}]
For each CCD the summary gives the $\chi^2/$ndf of the dark-current fit, and the gain, noise and dark current with
errors scaled by $\sqrt{\chi^2/\mathrm{ndf}}$. \texttt{*} marks a gain held at the all-peaks fit. The \emph{DC
bias} is the fitted dark current minus the one from the 0\,e$^-$/1\,e$^-$ pixel counts; it should be close to 0.
Below the table come the reason for each gain choice, the skipped CCDs, and the results of the optional steps.
\end{note}

%==================================================================================================================
\section{Gain versus charge: \texttt{-{}-gain-correction} and \texttt{gain\_vs\_charge.py}}\label{sec:gain}
%==================================================================================================================
The gain falls slowly with the charge in a pixel ($\sim$3\,\% at 1000\,e$^-$). Converting every pixel with the
low-charge gain therefore puts keV clusters low. \cc{gain_vs_charge.py} measures $g(n)$ from the spacing of the
electron comb at high charge, in the images themselves, using clean rows. It then converts each pixel's ADU to
electrons with that curve.

\subsection{Inside \texttt{run\_WADERs}}
\begin{center}
\begin{tabularx}{\textwidth}{@{}P{4.6cm}L@{}}
\toprule
\textbf{option} & \textbf{effect} \\
\midrule
\cc{--gain-correction} / \cc{pooled} & one curve per module from all its images $\to$ \cc{root_gaincorr/} \\
\cc{--gain-correction per-image} & each image's own curve; a CCD whose comb is unusable gets the module curve
  (noted in the summary) $\to$ \cc{root_gaincorr_per_image/} \\
\cc{--gain-correction both} & both, plus \cc{gain_vs_charge/per_image_vs_pooled.csv} (curves, energies,
  front-side K$\alpha$) \\
\cc{--gain-model} & form of $g(n)$: monotone interpolation through the measured points (default), two lines with
  a fitted knot, or one line \\
\cc{--gain-combine} & several images of a module: \emph{coherent} (default; each image converted with its own
  curve, then the pixels pooled, so the result is independent of the image order) or \emph{incoherent} (comb
  powers added) \\
\cc{--gain-anchor} & \emph{comb} (default; the comb's $g(n)$ throughout, which removes the image-to-image
  scatter of the fitted gains) or \emph{fit} (each image's fitted gain kept at low charge) \\
\bottomrule
\end{tabularx}
\end{center}
The plots in \cc{gain_vs_charge/} are \cc{<name>_gain_vs_charge.png} (local gain against charge) and
\cc{<name>_nonlinearity.png} ($A(n)/g(0)-n$ with its quadratic fit). CCDs with an untrustworthy curve (measured
to $<400$\,e$^-$, or rising) are copied unchanged.

\subsection{On its own}
\begin{lstlisting}
python src/gain_vs_charge.py measure img1.fz img2.fz ... -o OUTDIR [--name NAME] [--ext 1 3 4]
        [--model interpolated|broken|linear] [--combine coherent|incoherent] [--extrap-factor 3]
python src/gain_vs_charge.py apply OUTDIR/NAME_gain_vs_charge.json root/panaSKImg_clustersRec_*_ext*.root
        -o root_gaincorr [--anchor comb|fit]
\end{lstlisting}
\cc{measure} writes \cc{<name>_gain_vs_charge.json} (one model per extension) and the plots. \cc{apply}
rewrites \cc{pixels_E} and \cc{Energy} of each cluster file and copies \cc{info} and \cc{process_config}.
Method note: \cc{docs/gain_correction_method/gain_correction_method.tex}. Effect on the data:
\cc{docs/gain_correction_report.pdf} and \cc{docs/gain_correction_Am241.pdf}.

%==================================================================================================================
\section{Fe-55 composite analysis: \texttt{-{}-composite}}\label{sec:composite}
%==================================================================================================================
After the cluster files are written, \cc{MakeComposite_Fe55_Clusters_Extensions.cc} runs with ROOT, once per
image type and module. It produces the front/back composite clusters, the Fe-55 K$\alpha$/K$\beta$ fits and the
CTI maps.
\begin{itemize}
  \item \textbf{Input files}: \cc{root_gaincorr/} if \cc{--gain-correction} is on, else \cc{root/}. With
        \cc{--gain-correction both}, \cc{root_gaincorr_per_image/} is also used, giving
        \cc{composite_per_image/}. With \cc{--clusterer agmx}, the A\_GMX files are used instead; with
        \cc{both}, they are used in addition and the results go to \cc{composite_agmx/}.
  \item \textbf{Image type} $N$ (the CCD geometry in the macro: 4, 5, 6, 7, 8, 11): from \cc{Image_<N>} in the
        name, or \cc{--image-number}.
  \item \textbf{Module name}: from \cc{--module ACM=NAME}, else from a data folder named
        \cc{..._DMa_DMb_DMc_...} (ACM 110 $\to$ DMa, 109 $\to$ DMb, 106 $\to$ DMc), else \cc{ACM<id>}.
  \item \textbf{Output}: \cc{composite/Fe_55_Analysis_Plots/<module>/Image_<N>_<module>_list/} (\cc{ext1..4/}
        per CCD, \cc{compare/} 2$\times$2 overviews), \cc{composite/lists/}, and
        \cc{logs/composite_Image_<N>_<module>.log}.
\end{itemize}

%==================================================================================================================
\section{A\_GMX clustering}\label{sec:agmx}
%==================================================================================================================
A\_GMX (an adaptive regime switch plus Gaussian-mixture deblending) is the recommended alternative to panaSKImg's
ClusterFinder; its algorithm is described in \cc{docs/aGMX/aGMX_clustering.tex}. On 878 image-CCDs it gave:
\begin{itemize}
  \item unbiased Fe-55 energies at every cluster width;
  \item about 2.7$\times$ better energy resolution;
  \item resolved close X-ray pairs;
  \item few-electron events without fakes;
  \item a correct energy scale on crowded CCDs.
\end{itemize}
It runs on the calibrated images that panaSKImg writes (\cc{*_EXT<n>_waders.fits.bz2}) and produces cluster files
\textbf{in the WADERs format}, so every tool that reads WADERs cluster files reads them as well.

\subsection{Three ways to run it}
\begin{lstlisting}
# 1. in the pipeline (after panaSKImg and the gain correction)
run_WADERs images.txt --gain-correction --clusterer agmx --composite     # composite on A_GMX
run_WADERs images.txt --gain-correction --clusterer both --composite     # composite/ and composite_agmx/

# 2. on a finished run_WADERs folder
run_AGMX WADERs_analysis_<run> [-j 4] [--no-gaincorr] [--per-image] [--anchor comb|fit] [-o DIR]

# 3. the module directly
python src/agmx_root.py WADERs_analysis_<run> [-o OUTDIR] [--curve gain_vs_charge/ | curve.json] [-j N]
\end{lstlisting}

\begin{center}
\begin{tabularx}{\textwidth}{@{}P{4.6cm}L@{}}
\toprule
\textbf{output folder} & \textbf{clusters of} \\
\midrule
\cc{root_agmx/} & the images as calibrated by panaSKImg (fitted gain) \\
\cc{root_agmx_gaincorr/} & the images converted with the module's $g(n)$ curve from \cc{gain_vs_charge/} (same
  \cc{--gain-anchor} as \cc{root_gaincorr/}); written when the run has \cc{gain_vs_charge/} \\
\cc{root_agmx_gaincorr_per_image/} & the same with each image's own curve (\cc{run_AGMX --per-image}, or
  \cc{--gain-correction per-image/both}) \\
\bottomrule
\end{tabularx}
\end{center}
Each folder has \cc{root_files.txt}. Speed: about 30--60\,s per CCD (a crowded CCD takes longer). panaSKImg's
ClusterFinder still runs, because its file for each CCD supplies the \cc{info} and \cc{process_config} trees and
the seed settings. \cc{root/} is never changed.

\subsection{File format}
The A\_GMX clusters are turned into pysimdamicm \cc{Cluster} objects, the class the ClusterFinder itself uses.
Every branch is therefore computed by panaSKImg's own code, and the file is written with pysimdamicm's
\cc{OutputDataManager}:
\begin{itemize}
  \item \textbf{Same as WADERs}: trees \cc{clustersRec}, \cc{info} and \cc{process_config} (the last two copied
        from the CCD's panaSKImg file), and all WADERs branches: \cc{Energy}, \cc{PosX/Y}, \cc{wSTD_*}, \cc{fSTD_*},
        \cc{pixels_x/y/E}, \cc{Npix}, \cc{Qmax}, \cc{has_seed}, the PCA and track variables, \cc{closest_cls_*},
        \cc{Energy_PS/AVG}, \cc{RUNID}, \cc{ccd}, \cc{nfile} and \cc{readout_*}.
  \item \textbf{Same conventions}: \cc{pixels_x/y} are 1-based; \cc{pixels_E} is in keV with $e_{2\mathrm{eV}}$
        from \cc{process_config}, and \cc{Energy} $=\sum$\cc{pixels_E}. A pixel shared by two clusters is split
        between them, so the sum is the A\_GMX energy.
  \item \textbf{Extra branches} (one value per cluster):
\end{itemize}
\begin{center}
\begin{tabularx}{\textwidth}{@{}P{3.2cm}L@{}}
\toprule
\cc{agmx_cls} & class: 1 x-ray, 2 x-ray diffuse, 3 low-E, 4 muon, 5 worm, 6 alpha, 7 big deposit, 8/9 row/column
  artefact, 10 extended \\
\cc{agmx_blend} & fraction of the energy shared with neighbours (deblending) \\
\cc{agmx_ped} & local pedestal: 25th percentile of the charge 5--8\,px around the cluster [e$^-$/pix] \\
\cc{agmx_peakfrac} & largest Gaussian component / cluster \\
\cc{agmx_chi2} & $\chi^2$ per pixel of the cluster's own model \\
\cc{agmx_sig} & significance of the cluster \\
\cc{agmx_regime} & 0 sparse, 1 crowded, 2 unmodelled (hysteresis) \\
\bottomrule
\end{tabularx}
\end{center}
A flag value of $-999$ means ``not computed'': \cc{agmx_chi2} of clusters the mixture did not fit (few-electron
finder, hysteresis), or \cc{agmx_ped} with fewer than 15 active pixels in the ring (next to the mask). The file also
contains a \cc{TNamed agmx_config} with the version, the parameters and the codes, as JSON.

\begin{note}[Recommended selection for X-ray lines]
\cc{agmx_blend < 0.02} and \cc{agmx_ped < 10}, plus the analysis' usual $\sigma_{xy}$ cut (e.g.\ $<0.7$ for the
Am-241 lines, $<1$ for front-side Fe-55). The \cc{ped} cut removes clusters on a charge sheet. This is what made
the high-side tail of the 59.5\,keV line on the crowded CCD (docs/aGMX, section ``Contamination'').
\end{note}

\begin{note}[Checked against WADERs]
\cc{clustering_tests/agmx_root_check.py PANASKIMG_ROOT WADERS_FITS EXT} feeds the clusters of a real panaSKImg
file through the writer and compares every branch. On EXT1, EXT3 and EXT4 of an Am-241 image (1652, 7686 and
5452 clusters), every branch is identical except in numerically degenerate cases, where panaSKImg's own result
depends on the last bits of its inputs:
\begin{itemize}
  \item \cc{fSTD} of profiles 1--3 bins wide;
  \item the sign of PCA eigenvectors ($\theta$ versus $\theta\pm\pi$);
  \item 1--2 pixel PCA;
  \item \cc{closest_cls_*} when the nearest cluster lies in the same row or column. panaSKImg's ``exclude
        itself'' test is \cc{PosX != ... and PosY != ...}, and the writer uses the same rule.
\end{itemize}
\end{note}

\subsection{From Python}
\begin{lstlisting}[style=py]
import sys; sys.path.insert(0, "comm_scripts/clustering_tests")
import adaptive as AD
clusters = AD.adaptive_cluster(Q, valid, engine="gmx")     # Q: image in electrons, valid: bool mask
good = [c for c in clusters if c["blend"] < 0.02 and c.get("ped", 0) < 10 and not c.get("edge")]
# each cluster: E [keV], y, x, sxy [px], cls, blend, ped, peakfrac, chi2, sig, edge, regime
# return_pixels=True also gives c["_pix"] = (rows, cols, electrons), summing exactly to E
\end{lstlisting}
\cc{src/agmx_root.py} exposes the steps separately: \cc{run(run_dir, outdir, curve, anchor, jobs)},
\cc{process(...)} for one CCD, \cc{build_clusters(...)}, \cc{write_root(...)} and \cc{read_reference(...)}.

%==================================================================================================================
\section{The pipeline modules on their own}\label{sec:modules}
%==================================================================================================================
\subsection{\texttt{autoconfig\_panaSKImg.py}: a panaSKImg JSON from a raw image}
\begin{lstlisting}
python src/autoconfig_panaSKImg.py image.fz -o image.json [options]
python src/autoconfig_panaSKImg.py image.fz -o OUTDIR --per-ext      # one JSON per CCD + run_<image>.sh
\end{lstlisting}
It measures geometry, gain, noise, dark current, dead channels, bad columns and the pedestal method from the raw
pixels. It changes only these settings in the reference JSON, which is built in (\cc{-r my.json} uses your own).
The same options can be passed through \cc{run_WADERs --autoconfig-args "..."}.

\begin{longtable}{@{}P{5.2cm}P{10.8cm}@{}}
\caption{autoconfig options (also usable inside \texttt{-{}-autoconfig-args}).}\label{tab:autoconfig}\\
\toprule
\textbf{option} & \textbf{meaning (default)} \\
\midrule
\endhead
\cc{--gain G} & gain prior [ADU/e$^-$] if peaks are not resolved; also per CCD \cc{EXT1=1.32,EXT2=1.30} \\
\cc{--gain-min}, \cc{--gain-max} & gain search range (0.6, 3.0) \\
\cc{--prior-tol} & allowed relative difference measured gain vs prior (0.25) \\
\cc{--max-sig-over-gain} & noise/gain above which the peaks count as unresolved (0.3) \\
\cc{--lam-max L} & rows with more dark current than $L$ e$^-$/pix are left out of the fit (1.0) \\
\cc{--min-rows N} & minimum number of rows in the fit (25) \\
\cc{--peaks auto|all|N} & peaks to fit: \cc{auto} = those that follow one Poisson distribution (default) \\
\cc{--poisson-tol}, \cc{--max-peaks} & tolerance (0.25) and upper limit (10) of \cc{--peaks auto} \\
\cc{--no-fallback} & no held-gain 0\,e$^-$/1\,e$^-$ fit for unresolved CCDs \\
\cc{--ccd-cols N} & prescan+active columns (default from the header) \\
\cc{--dead-frac F}, \cc{--keep-dead} & detection of channels without a CCD ($2\times10^{-4}$), or keep them \\
\cc{--cluster-threshold-e P S | auto} & ClusterFinder thresholds in electrons (pixel $>P$, seed $>S$), or
  \cc{auto} per CCD from its dark charge (smallest $k$ with $P(\mathrm{Poisson}(\lambda)\ge k)\le10\,\%$, pixel
  $>k-0.5$, seed $>9.5$\,e$^-$). Default: the reference's 10$\sigma$ / 30$\sigma$ \\
\cc{--cluster-all} & also cluster flooded/outlier CCDs \\
\cc{--plot-peaks N} & DC-fit plot shown up to $N$ e$^-$ (5; needs patched pysimdamicm, see
  \cc{docs/patches/}) \\
\cc{--tune}, \cc{--tune-dir} & choose the pedestal method/window with panaSKImg; keep those runs \\
\cc{--per-ext} & \cc{-o} is a folder: one JSON per CCD and a run script \\
\cc{--panaskimg PATH} & panaSKImg used by \cc{--tune} and the run script \\
\bottomrule
\end{longtable}
In \cc{--autoconfig-args}, do not pass \cc{--peaks}, \cc{--per-ext} or \cc{--tune}; \cc{run_WADERs} sets these
itself.

\subsection{\texttt{gain\_dc\_split.py}: gain from all peaks, dark current from two}
\begin{lstlisting}
python src/gain_dc_split.py image.fz -o OUTDIR [--per-ext] [--run-full OUT] [-j 4] [--tune]
        [--autoconfig-args "..."] [--panaskimg PATH]
\end{lstlisting}
This is step 1 of \cc{run_WADERs}. It writes \cc{gain/}, \cc{dc/} and \cc{final/} JSONs and a summary comparing
the fits; \cc{--run-full OUT} also runs the full panaSKImg sequence.

\subsection{\texttt{fit\_quality.py}: how good are the dark-current fits}
\begin{lstlisting}
python src/fit_quality.py FOLDER [--ref FOLDER ...] [--csv out.csv]
\end{lstlisting}
FOLDER is a panaSKImg output folder or a \cc{run_WADERs} output folder (fits in \cc{root/dc_fits/}, references in
\cc{json/<image>/}). It reports $\chi^2/$ndf, the scaled errors, and the dark-current bias against the
0\,e$^-$/1\,e$^-$ counts of the autoconfig \cc{*_measured.json}.

\subsection{\texttt{gain\_vs\_charge.py} and \texttt{agmx\_root.py}}
See Sections~\ref{sec:gain} and~\ref{sec:agmx}.

%==================================================================================================================
\section{Utilities}\label{sec:utils}
%==================================================================================================================
\begin{longtable}{@{}P{4.6cm}P{11.4cm}@{}}
\toprule
\textbf{script} & \textbf{what it does and how to run it} \\
\midrule
\endhead
\cc{compare_runs.py} & Are two \cc{run_WADERs} outputs identical? It compares \cc{summary.txt}, \cc{fit_quality.csv},
  the final JSONs, the $g(n)$ curves, every cluster energy in \cc{root*/} and the composite tables; a file present
  in only one run is a difference. PDFs and logs are not compared. Exit code 0 if identical. \newline \cc{python utils/compare_runs.py OLD NEW [-q]} \\
\cc{split_clustersRec_by_ext.py} & Splits files made with \cc{--acm} (all CCDs stacked, \cc{_ext1234.root}) into
  one file per CCD. \newline \cc{python utils/split_clustersRec_by_ext.py files... -o OUTDIR [--nrows 500] [--list L]} \\
\cc{compare_dcfit_vs_skips.py} & Tiles the per-CCD plots of runs with different skip numbers
  (\cc{analysis/<skips>/}) into one PDF per CCD. \newline
  \cc{python utils/compare_dcfit_vs_skips.py analysis/ -o OUT [--ncols 4] [--plot FitDarkCurrentProcess_DCfit]} \\
\cc{skip_noise_L1.py} & Decodes a level-1 image (every skip kept, \cc{*.RJ45ADCDATA<ch>}) and plots noise against
  the number of skips, the noise spectrum and the start-of-pixel transient. \cc{--composites N} also writes $N$
  partial averages. Outputs are named after the file plus \cc{_ch<k>}. It needs \cc{src/}. \newline
  \cc{python utils/skip_noise_L1.py FILE -o OUT [--geometry RxCxS] [--skip-first 20] [--gain G] [--max-rows N]} \\
\bottomrule
\end{longtable}

%==================================================================================================================
\section{Research code: \texttt{clustering\_tests/}}\label{sec:clustering}
%==================================================================================================================
This folder holds every clustering method compared with WADERs in the study of 2026-10-23 to 2026-10-25, together
with the scripts that tested them. Apart from A\_GMX, which the writer uses, none of it is part of the pipeline.
\begin{itemize}
  \item \textbf{Paths}: the scripts read data from absolute paths under \cc{~/Documents/}.
  \item \textbf{Outputs}: written next to the scripts (git-ignored).
  \item \textbf{Saved results}:
        \cc{~/Documents/comm_scripts_data/session_2026-10-23_clustering/clustering_2026-10-25/}.
  \item \textbf{Python}: the \cc{damic} environment.
\end{itemize}

\subsection{Methods}
\begin{center}
\begin{tabularx}{\textwidth}{@{}P{3.3cm}P{2.4cm}L@{}}
\toprule
\textbf{module} & \textbf{method} & \textbf{entry point and idea} \\
\midrule
\cc{adaptive.py} & \textbf{A\_GMX} & \cc{adaptive_cluster(Q, valid, engine="gmx")}: regions the model cannot
  describe go to hysteresis; everything else goes to GMX \\
\cc{gmx.py} & GMX & \cc{general_cluster(Q, valid, bg)}: Gaussian-mixture EM deblending, ridge/shape merging,
  island BIC, quality flags \\
\cc{mhc.py} & MHC & specialised finders (X-ray, track, streak, low-E, \dots) with $\chi^2$/BIC arbitration and
  classes \\
\cc{methods2.py} & hybrid v1/v2 & MFD+G pieces merged by B-region and saddle tests (\cc{mfdg_hybrid}),
  track-first \\
\cc{agent_methods.py} & MFD+G & background map, Poisson tails, matched-filter detection + disk + growth \\
\cc{cluster_methods.py} & WADERs-like, B & hysteresis (\cc{wader}), topological clustering in electrons
  (\cc{topo}), injection (\cc{inject}) \\
\cc{pws.py} & PWS & persistence-simplified watershed (a control; not recommended) \\
\bottomrule
\end{tabularx}
\end{center}

\subsection{Test and validation scripts}
\begin{longtable}{@{}P{6.3cm}P{9.7cm}@{}}
\toprule
\textbf{script(s)} & \textbf{test} \\
\midrule
\endhead
\cc{wide_test.py}, \cc{wide_run.py adaptive|GMX|adaptive+GMX} (+ \cc{summ_wide.py}) &
  every method on all 878 calibrated image-CCDs: injection, real K$\alpha$, 59.5\,keV \\
\cc{wide_agmx_flags.py} (+ \cc{summ_wide_flags.py}) & Fe-55 K$\alpha$ with the quality-flag cuts; injection \\
\cc{am241_lines_flags.py} (\cc{... fit}) & every Am-241 line, old vs new A\_GMX, with/without the flag cuts \\
\cc{adaptive_validate.py} (+ \cc{run_*.sh}) & tracks, few-electron, Fe-55 and crowded-CCD validation \\
\cc{compare.py}, \cc{mhc_fe55.py} (+ \cc{summ_fe55.py}) & Fe-55 injection at 4 widths, pairs, real K$\alpha$ \\
\cc{am_test.py} & Am-241 lines, injected 17.75 / 59.5\,keV, splitting of big events \\
\cc{tracks_test.py}, \cc{extscan.py}, \cc{mhc_tracks.py} & injected muons, worms, blobs, alphas, hot columns,
  pairs \\
\cc{lowE_test.py}, \cc{mhc_lowE.py} & few-electron events: efficiency, exact charge, fakes \\
\cc{am241_29img_recluster.py}, \cc{am241_29img_fit.py}, \cc{adaptive_*_am241.py}, \cc{make_crowded_table.py} &
  crowded Am-241 CCD: the constrained Lepy fit for each clustering; occupancy and threshold scans \\
\cc{agmx_root_check.py} & the WADERs-format check of \cc{src/agmx_root.py} (Section~\ref{sec:agmx}) \\
\bottomrule
\end{longtable}

\subsection{Documentation builders}
\begin{center}
\begin{tabularx}{\textwidth}{@{}P{6.3cm}L@{}}
\toprule
\cc{make_agmx_figures.py} & \cc{docs/aGMX/figures/fig00-13}: every A\_GMX step on a real image \\
\cc{make_agmx_perf.py} & \cc{fig20-26} and \cc{docs/aGMX/perf_tables.tex} \\
\cc{make_agmx_contamination_figs.py} & \cc{fig30-34}: the 59.5\,keV tail and the quality flags \\
\cc{make_overleaf_zip.py} & \cc{docs/aGMX/aGMX_clustering_overleaf.zip} \\
\cc{make_clustering_report.py} (+ \cc{report_*.py}) & \cc{docs/clustering_methods_report.pdf} \\
\bottomrule
\end{tabularx}
\end{center}

%==================================================================================================================
\section{Analysis notes and reports}\label{sec:notes}
%==================================================================================================================
\subsection{Am-241 lines, 10--30 keV (\texttt{docs/am241\_note/})}
This note covers the gain correction per gain period, the cluster-width cut $\sigma_{xy}<0.7$, the constrained fit
of the Np L lines (Lepy et al.\ ratios, one energy scale $k$), and the two extra lines at 22.1 and 24.9\,keV. Its
input is two \cc{run_WADERs} outputs of the same images (\cc{no_gaincorr/} and \cc{gaincorr/}). Defaults are
\cc{--run-dir ~/Documents/JH1_Testing/AMTEst/runs_29img} (or \cc{$AM241_RUN_DIR}) and \cc{--run-name AMtest_29}.
\begin{center}
\begin{tabularx}{\textwidth}{@{}cp{5.6cm}L@{}}
\toprule
\textbf{step} & \textbf{run} & \textbf{does} \\
\midrule
0 & \cc{run_WADERs list.txt -o no_gaincorr --composite} \newline \cc{run_WADERs list.txt -o gaincorr --composite --gain-correction} &
  the pipeline ($\sim$35 min each for 29 images) \\
1 & \cc{correct_gain_per_period.py --list list.txt --periods 21-40 41-50} & $g(n)$ per gain period \\
2 & \cc{extract_clusters.py} & cluster files $\to$ one table per CCD \\
3 & \cc{fit_all_peaks.py [--no-grid]} & all fits; the band-cut grid ($\sim$40 min) \\
4 & \cc{peak_table.py} & mean, $\sigma$, FWHM and counts of every peak \\
5 & \cc{identification_tests.py} & tests of the extra-line hypotheses \\
6 & \cc{make_note_figures.py} & every figure and table of the note \\
-- & \cc{gold_check.py}, \cc{make_report.py} & gold L lines; the companion report \\
\bottomrule
\end{tabularx}
\end{center}
The LaTeX source is \cc{note/main.tex}; the Overleaf upload is \cc{docs/AMtest_Am241_10-30keV_note_overleaf.zip}.

\subsection{Gain correction}
\begin{itemize}
  \item \cc{docs/gain_correction_method/gain_correction_method.tex}: the method note. Its figures are made by
        \cc{make_figures.py}, and \cc{am241_forms.py} compares the three model forms on Am-241.
  \item \cc{docs/report_scripts/make_gain_correction_report.py} $\to$ \cc{docs/gain_correction_report.pdf}
        (Fe-55 across the production data). \cc{make_gain_correction_Am241.py} $\to$
        \cc{docs/gain_correction_Am241.pdf}.
\end{itemize}

\subsection{Clustering}
\begin{itemize}
  \item \cc{docs/aGMX/aGMX_clustering.tex} (Overleaf zip next to it): A\_GMX step by step, with equations,
        figures, performance against WADERs, and the contamination study with its quality flags.
  \item \cc{docs/clustering_methods_report.pdf}: all methods and all tests.
\end{itemize}

\subsection{Hand-over notes}
\cc{docs/README.md} is the index of the documentation; \cc{docs/PIPELINE_GUIDE.md} gives the physics and the code
map. \cc{docs/notes/} holds the dated hand-over notes (each starts with a status line: current, or what supersedes
it); \cc{CONTEXT_2026-10-24.md} is the current one, and \cc{OPEN_ITEMS.md} lists everything still open in one place.

%==================================================================================================================
\section{Tests and changing the code safely}\label{sec:tests}
%==================================================================================================================
\begin{lstlisting}
cd comm_scripts
python -m unittest discover -s tests -v          # with the damic / panaSKImg python
\end{lstlisting}
\begin{center}
\begin{tabularx}{\textwidth}{@{}P{5.3cm}L@{}}
\toprule
\cc{test_run_wader.py} & input handling of \cc{run_WADERs} (lists, names, options) \\
\cc{test_gain_vs_charge.py} & the $g(n)$ models and their inversion \\
\cc{test_am241_model.py} & the constrained Am-241 fit on synthetic spectra \\
\cc{test_cluster_thresholds.py} & the automatic cluster thresholds in electrons \\
\cc{test_clustering_methods.py} & MFD+G, hybrid, MHC and the GMX quality flags on a simulated image \\
\cc{test_agmx_root.py} & the A\_GMX writer: file structure, panaSKImg conventions, flags, a reference without
  clusters, binned columns (simulated CCD) \\
\cc{test_utils.py} & \cc{compare_runs.py} on fake runs (relative paths, extra files, changed JSON) \\
\bottomrule
\end{tabularx}
\end{center}
For a change to the pipeline itself, run \cc{run_WADERs} on the same images with the old and the new code, into
different \cc{-o} folders, and compare them with \cc{utils/compare_runs.py OLD NEW}. The pipeline is
deterministic, so any difference comes from the change. A git worktree of the old commit is a convenient way to
run the old code.

%==================================================================================================================
\section{Troubleshooting}
%==================================================================================================================
\begin{longtable}{@{}P{5.6cm}P{10.4cm}@{}}
\toprule
\textbf{symptom} & \textbf{cause / fix} \\
\midrule
\endhead
``already exists; use \cc{--overwrite}'' & the output folder is not empty: \cc{--overwrite}, or another \cc{-o} \\
``no conda environment with panaSKImg'' & \cc{--env NAME} / \cc{$WADERS_ENV}, or activate the environment \\
a CCD is ``not run (flooded/unresolved)'' & the peaks could not be resolved; \cc{--all-ext} forces it (fits not
  reliable) \\
very slow clustering & high dark current: \cc{--autoconfig-args "--cluster-threshold-e 2 5"} (or \cc{auto}) \\
MakeComposite: ``not run: no Image\_<N>'' & give \cc{--image-number N} \\
composite module called ACM<id> & give \cc{--module <id>=<name>} \\
``no reliable g(n) curve'' (A\_GMX log) & that CCD is clustered with its fitted gain, as \cc{root_gaincorr/} keeps it unchanged \\
\cc{x_max_plot} not supported & unpatched pysimdamicm: the wider DC-fit plot is left out
  (\cc{docs/patches/pysimdamicm_x_max_plot.patch}) \\
\cc{_ext} in the output path (warning) & MakeComposite reads the extension from \cc{_ext<N>}; avoid it in folder names \\
\bottomrule
\end{longtable}

\end{document}
